\documentclass[12pt]{exam}

\usepackage{mystyle}
\usepackage{listings}
\usepackage{xcolor}

\definecolor{codegreen}{rgb}{0,0.6,0}
\definecolor{codegray}{rgb}{0.5,0.5,0.5}
\definecolor{codepurple}{rgb}{0.58,0,0.82}
\definecolor{backcolour}{rgb}{0.95,0.95,0.92}

\lstdefinestyle{mystyle}{
    backgroundcolor=\color{backcolour},
    commentstyle=\color{codegreen},
    keywordstyle=\color{blue},
    numberstyle=\color{codegray},
    stringstyle=\color{codepurple},
    basicstyle=\ttfamily\footnotesize,
    breakatwhitespace=false,
    breaklines=true,
    captionpos=b,
    keepspaces=true,
    numbers=left,
    numbersep=5pt,
    showspaces=false,
    showstringspaces=false,
    showtabs=false,
    tabsize=2
}
\lstset{style=mystyle}

\title{MATH 3191 Final Exam Review}

\date{}

\begin{document}
\maketitle
\section{Exam Format}
The exam will take the following format
\begin{enumerate}
    \item Some short response questions that require a sentence or two explaining why a statement
        is true or false (A beginning of what a proof might look like)
    \item Longer response questions that will require you to compute a numerical answer and show work
        on how it was achieved
    \item Questions on lecture based content (No video exclusive content, but from the pdfs themselves)
    \item A separate jupyter notebook document that I want you to fix (think of it like a really
        short python lab with the same content as previous ones)
\end{enumerate}
Note: Since this is a take-home exam, this exam will be fairly long (4-6 questions in each section.
I haven't made it yet, so this is a rough estimate)
I will time this as if you have 4 hours to take it (1 more hour than the last 2!), but you have the
entire week, so use your time wisely! There will be a focus on material after exam 2, but the
following content from previous exams are also fair game
\begin{enumerate}
    \item Diagonalization
    \item Similarity Transformations
    \item Fundamental Subspaces ($\colS{A}, \nullS{A}, \rowS{A}, \nullS{A^\top}$)
    \item Linear Independence and Dependence
    \item Computing a Basis or showing something is a basis
    \item Orthogonal Projections
    \item Orthogonality
    \item Coordinate Vectors (our $\mathcal{B}$ coordinates!)
\end{enumerate}

\section{Written Homework Problems}
You should be able to do all the problems from HW1 -- HW12. Filter HW1 -- HW11 via the above list.
\section{MyOpenMath Problems}
I would pay attention to the following topics associated with each problem below
\begin{enumerate}
    \item HW1 -- HW11: See previous review documents and filter by above list
    \item HW12: 3, 7, 9, 12, 16
        \begin{enumerate}
            \item Determine if a set is an orthogonal set (with respect to our standard inner product)
            \item Apply Gram-Schmidt to a list of vectors
            \item Compute a QR decomposition of an $m \times n$ matrix where $m, n \leq 3$
            \item Compute a least squares solution
            \item Compute an orthogonal projection of a vector onto a column space of a vector
        \end{enumerate}
    \item HW13: 2, 5, 7, 8, 11, 17, 18, 19, 23
        \begin{enumerate}
            \item Interpret the meaning of a slope with respect to a least squares solution
            \item Solve a regression problem with at most 3 variables $(z = \beta_0 + \beta_1 x + \beta_2 y)$
                or $(z = \beta_0 + \beta_1 x + \beta_2 x^2)$ given the data in either tabular or $(x,y)$ format.
            \item Predict a value from a regression equation
            \item Compute an inner product of given input vectors
            \item Determine if a matrix is orthogonal
            \item Compute a value to ensure a given matrix is orthogonal
            \item Orthogonally diagonalize a real symmetric matrix
        \end{enumerate}
\end{enumerate}
\section{Lecture Based Content}
See the previous review document filtering by the list at the top of this document.
\begin{enumerate}
    \item Explain how the SVD relates to an Eigendecomposition
    \item Apply Gram-Schmidt to vectors with a given inner product (Not our standard one!)
    \item Explain Python code for topics related to least squares, qr, and SVD
    \item Transform a ``weird'' space into $\R^n$ and perform some kind of decomposition with it
        (IE diagonalize a matrix of polynomials, and interpret the output in terms of polynomials)
\end{enumerate}
\end{document}
